% =============================================================================
\section{Buyback-and-Burn Mechanism}
\label{sec:buyback-burn}
% =============================================================================

This section specifies the Assistance Fund mechanism, describes the revenue
sources feeding continuous buyback-and-burn, models the circulating supply
trajectory, and proves conditions under which the supply becomes net
deflationary.

\subsection{The Assistance Fund}
\label{subsec:assistance-fund}

The Assistance Fund is the central value-accrual engine of PURECORTEX. It is
an on-chain smart contract account (implemented as an Algorand LogicSig) that:

\begin{enumerate}
  \item Receives \textbf{90\% of all protocol revenue} from every revenue
    source (Table~\ref{tab:revenue-sources}).
  \item \textbf{Continuously purchases CORTEX} from the primary CORTEX/ALGO
    DEX liquidity pool.
  \item \textbf{Immediately burns} all purchased tokens by sending them to the
    Algorand zero address (an irrecoverable sink) in the same atomic
    transaction group.
\end{enumerate}

The Assistance Fund is seeded with 5\% of the total CORTEX supply at genesis,
providing initial operational capital for the buyback engine. After launch, it
is entirely self-sustaining: the 90\% revenue stream ensures continuous
buyback-and-burn activity proportional to protocol usage.

This model is inspired by the mechanism employed by the leading decentralized
perpetuals protocol, which directs the vast majority of its protocol revenue to
an analogous buyback fund, creating sustained deflationary pressure that
directly rewards token holders through supply reduction rather than dividend
distribution.

\subsection{Revenue Sources}
\label{subsec:revenue-sources}

All protocol revenue flows through the 90/10 split defined in
Table~\ref{tab:fee-split} (Section~\ref{sec:dynamic-fees}). The revenue
sources feeding the Assistance Fund are:

\begin{table}[H]
\centering
\caption{Revenue sources for the Assistance Fund.}
\label{tab:revenue-sources}
\begin{tabular}{@{}llr@{}}
\toprule
\textbf{Source} & \textbf{Description} & \textbf{To Assistance Fund} \\
\midrule
Trading Fees & Dynamic 0.5\%--2.0\% on all bonding curve trades & 90\% \\
Agent Creation Fee & 100 CORTEX per agent created & 90\% \\
x402 Commission & 5\% tax on inter-agent MCP micropayments & 90\% \\
LP Fees & Protocol-owned LP position accrued fees & 90\% \\
\midrule
\multicolumn{2}{@{}l}{\textbf{Remaining 10\% of all sources}} & Operations \\
\bottomrule
\end{tabular}
\end{table}

\subsection{Continuous Buyback Mechanism}
\label{subsec:buyback-mechanism}

Unlike capped weekly buyback models, the Assistance Fund operates
\emph{continuously}: revenue is converted to buyback-and-burn operations as
it accumulates, with no weekly caps or rate limits.

\begin{definition}[Continuous Buyback]
\label{def:continuous-buyback}
Let $\mathcal{R}(t)$ denote the cumulative protocol revenue at time $t$. The
cumulative CORTEX purchased and burned by the Assistance Fund is:
\begin{equation}
  B(t) = \int_0^t \frac{0.90 \cdot r(\tau)}{P(\tau)} \, d\tau + B_{\text{seed}}
  \label{eq:continuous-buyback}
\end{equation}
where $r(\tau)$ is the instantaneous revenue rate, $P(\tau)$ is the CORTEX
market price at time $\tau$, and $B_{\text{seed}} = 0.05 \cdot S_{\text{total}}$
is the initial seed allocation available for burn.
\end{definition}

The absence of weekly caps is a deliberate design choice: caps protect against
market manipulation in treasury-discretion models, but the Assistance Fund is
fully automated with no discretionary spending authority. The burn rate is
mathematically bounded by protocol revenue, which itself is bounded by trading
activity --- providing a natural, market-driven rate limit.

Purchased tokens are immediately burned in the same atomic transaction group,
ensuring that bought-back tokens cannot re-enter circulation.

\subsection{Circulating Supply Model}
\label{subsec:supply-model}

The circulating supply at time $t$ is governed by:

\begin{equation}
  S(t) = S_0 + E(t) + V(t) - B(t)
  \label{eq:supply-model}
\end{equation}

where:
\begin{itemize}
  \item $S_0$ is the initial circulating supply at TGE (liquidity allocation,
    genesis distribution, and creator TGE release).
  \item $E(t) = \int_0^t e(\tau) \, d\tau$ is the cumulative staking emissions
    released by time $t$, with $e(\tau)$ the instantaneous emission rate from
    Equation~\eqref{eq:staking-emission}.
  \item $V(t)$ is the cumulative creator vesting releases.
  \item $B(t) = \int_0^t b(\tau) \, d\tau$ is the cumulative tokens burned by
    time $t$ via the Assistance Fund's continuous buyback-and-burn
    (Equation~\eqref{eq:continuous-buyback}).
\end{itemize}

\subsection{Deflationary Crossover}
\label{subsec:deflationary}

\begin{theorem}[Deflationary Crossover]
\label{thm:deflationary}
The circulating supply $S(t)$ is decreasing (net deflationary) at time $t$ if
and only if:
\begin{equation}
  \frac{dB}{dt}(t) > \frac{dE}{dt}(t) + \frac{dV}{dt}(t)
  \label{eq:deflation-condition}
\end{equation}

That is, the instantaneous burn rate exceeds the sum of the emission rate and
vesting release rate.
\end{theorem}

\begin{proof}
Differentiating Equation~\eqref{eq:supply-model}:
\begin{equation}
  \frac{dS}{dt} = \frac{dE}{dt} + \frac{dV}{dt} - \frac{dB}{dt}
  \label{eq:ds-dt}
\end{equation}

The supply is decreasing when $dS/dt < 0$, which rearranges to:
\begin{equation}
  \frac{dB}{dt} > \frac{dE}{dt} + \frac{dV}{dt}
  \label{eq:deflation-rearranged}
\end{equation}
\qed
\end{proof}

\subsection{Crossover Timing Analysis}
\label{subsec:crossover-timing}

We analyze the conditions under which Equation~\eqref{eq:deflation-condition}
is satisfied. The emission rate is determined by the schedule in
Equation~\eqref{eq:staking-emission}:

\begin{equation}
  \frac{dE}{dt} = \frac{E(y)}{365} \quad \text{for year } y
  \label{eq:emission-rate}
\end{equation}

In year 1, this equals $9.6 \times 10^{14} / 365 \approx 2.63 \times 10^{12}$
tokens per day. By year 3, it declines to $4.8 \times 10^{14} / 365 \approx
1.32 \times 10^{12}$ tokens per day.

The vesting release rate consists solely of the creator allocation (10\% of
supply, vesting over 180 days): $0.90 \times 1.0 \times 10^{15} / 180 \approx
5.0 \times 10^{12}$ per day during the first 180 days, and zero thereafter.
There is no team or advisor vesting.

The burn rate depends on protocol revenue, which grows with the number of
active agents and trading volume. With 90\% of all revenue flowing to the
Assistance Fund, the burn rate is:
\begin{equation}
  \frac{dB}{dt} \approx 0.90 \cdot A(t) \cdot \bar{V}(t) \cdot \bar{f} \cdot \frac{1}{\bar{P}(t)}
  \label{eq:burn-rate-approx}
\end{equation}

where $A(t)$ is the number of graduated agents, $\bar{V}(t)$ is the average
daily trading volume per agent, $\bar{f}$ is the average fee rate, and
$\bar{P}(t)$ is the average CORTEX price.

\begin{remark}
The 90\% revenue allocation to burns dramatically accelerates the deflationary
crossover compared to models where only a fraction of treasury revenue funds
buybacks. Under moderate growth assumptions (100 graduated agents by month 12,
500 by month 18, with average daily volume of 10{,}000 CORTEX per agent),
Monte Carlo simulations (Section~\ref{sec:simulations}) place the deflationary
crossover at approximately month 12--18. Under conservative assumptions
(50 agents by month 12), the crossover extends to month 18--24. Under
aggressive growth (200 agents by month 12), the crossover can occur as early
as month 8--12.
\end{remark}

\subsection{Long-Run Supply Convergence}
\label{subsec:long-run}

After the staking emission schedule concludes (year 4) and creator vesting
completes (day 180), the supply dynamics simplify to:

\begin{equation}
  \frac{dS}{dt} = -\frac{dB}{dt} < 0
  \label{eq:long-run}
\end{equation}

Under the assumption that protocol revenue remains positive (at least one
active agent generating trading volume), the supply is strictly decreasing for
all $t > 4$ years. The asymptotic supply converges to:

\begin{equation}
  S_\infty = \lim_{t \to \infty} S(t) = S(4) - \int_4^\infty b(\tau) \, d\tau
  \label{eq:s-infinity}
\end{equation}

If the burn rate is bounded below by some $b_{\min} > 0$, then $S_\infty = 0$
in finite time. In practice, the burn rate is proportional to protocol revenue,
which itself depends on the token economy's activity level. As supply
decreases, token scarcity increases price, which reduces the number of tokens
purchased per unit of revenue:
\begin{equation}
  \frac{dB}{dt} = \frac{0.90 \cdot r(t)}{P(t)} \to 0
  \quad \text{as } P(t) \to \infty
  \label{eq:equilibrium-supply}
\end{equation}

This creates a deflationary attractor: the supply continuously decreases but
never reaches zero, with the rate of decrease slowing as the per-token price
increases---a property analogous to exponential decay with a time-varying rate
constant. The absence of weekly caps means the burn rate responds immediately
to changes in protocol revenue, creating tight coupling between ecosystem
activity and deflationary pressure.
